\chapter{Introduction}
\label{sec:intro}

\noindent
Modern sequence generators are the engines behind a wide range of transformative applications. They translate sentences from one language to another in machine translation systems, convert audio into text in automatic speech recognition (ASR) systems, and produce natural language descriptions of visual scenes in image captioning. In all of these environments, an autoregressive model proposes many possible outputs token by token, and a decoding algorithm must decide which full sequence to present to the user. The historically dominant practice is to prefer the sequence that the model believes is most likely. Procedures such as greedy search and beam search operationalize this idea by approximating Maximum A Posteriori (MAP) decoding, effectively selecting the single hypothesis that maximizes the internal probability distribution of the model.

Over the last few years, however, evidence has mounted that raw probability is an imperfect stand-in for what humans actually value. In machine translation, sentences judged to be of higher quality by advanced neural evaluation metrics often receive a lower model probability than the generic winners selected by beam search. Similar mismatches appear in ASR and image captioning, where the highest-confidence transcript or description is not necessarily the most semantically accurate or visually grounded. Because model likelihood is strictly optimized for next-token prediction rather than for downstream utility, likelihood-first decisions frequently fall short of the true performance ceiling of the foundation model. These observations strongly motivate a paradigm shift from likelihood-first decoding to quality-first selection.

This work proposes the development and evaluation of Quality-Aware Decoding (QAD) strategies as robust alternatives for multimodal generative fields. The core premise is straightforward: instead of committing prematurely to a single high-probability output, the system generates a diverse pool of plausible candidates and subsequently selects the one that scores highest under a task-appropriate notion of quality. For ASR, quality is defined by minimizing transcription errors and preserving semantic intent, captured by metrics such as Word Error Rate (WER) and referenceless evaluators. For image captioning, quality encompasses visual relevance and structural consensus, measured by continuous embedding similarities and discrete $n$-gram metrics. Across both modalities, the shared architectural pattern is a generate-then-rerank pipeline driven by the specific utility that best reflects human judgment.

A principled mathematical framework for formalizing this shift is Minimum Bayes Risk (MBR) decoding. Rather than asking which output is most probable, MBR asks which output exhibits the best expected utility when compared against the distribution of plausible alternatives. In practice, this expectation is approximated by generating a diverse candidate list and performing pairwise evaluations to find the ultimate consensus sequence. While MBR consistently improves upon MAP decoding, this thesis identifies and addresses critical vulnerabilities in its deployment. Through extensive empirical evaluation, this research exposes a severe hallucination trap where text-centric estimators blindly reward fluent but acoustically disconnected outputs. Furthermore, it highlights a strict Pareto frontier in vision-language tasks, where maximizing pure visual grounding severely degrades linguistic structure. To circumvent these vulnerabilities and entirely bypass the quadratic computational bottleneck of traditional MBR, this work proposes and validates a Two-Stage MBR framework. By leveraging fast referenceless estimators for initial candidate pruning followed by rigorous semantic consensus, this hybrid approach achieves near-optimal quality while significantly reducing inference latency.

The remainder of this document is organized as follows. Chapter 2 introduces the foundational concepts utilized throughout the work, detailing sequence-to-sequence architectures, traditional decoding strategies, and the evaluation metrics defining task quality. Chapter 3 surveys the related work, contextualizing QAD, test-time compute, and the known limitations of MBR across various modalities. Chapter 4 formally presents the methodology adopted in this thesis, defining the mathematical mechanics of candidate generation, risk-based selection, and the proposed two-stage fusion. Chapter 5 outlines the experimental setup, detailing the hardware, dataset benchmarks, and model architectures utilized to evaluate the pipeline. Chapter 6 presents the empirical results and discussion, rigorously analyzing the impact of candidate diversity, the referenceless hallucination trap, and the computational efficiency of test-time scaling. Finally, Chapter 7 concludes the thesis, summarizing the primary achievements and identifying highly promising directions for future multimodal research.
